% ============================================
% Section: Conclusion
% ============================================
\label{sec:conclusion}

We presented Co-LIC2, a cooperative multi-node photo-realistic Gaussian
SLAM system whose central contribution is a purpose-built wireless
communication architecture for bandwidth-constrained, lossy, heterogeneous
mobile environments. The system contributes a QoS-aware multi-stream protocol
mapping four traffic classes to 802.11e EDCA queues with application-layer
reliability; a submap-aware adaptive compression scheduler jointly optimizing
SH-degree, quantization, and entropy-coding under a bandwidth budget using
the perceptually motivated NPV metric; a hierarchical hybrid network
architecture with stable cluster-head election, multi-hop submap routing,
and three graceful degradation modes; and a cooperative pose-graph back-end
for globally consistent merging of heterogeneous Gaussian maps.

The protocol is fully specified with finite-state machines, message
catalogues, bandwidth models, and failure-mode recovery analysis. Analytical
bounds demonstrate per-node uplink as low as 3.2~Mbps, per-cluster submap
merge latency under 2~s, and sub-linear per-link bandwidth growth with
node count. Co-LIC2 is implemented as a backward-compatible extension of the
open-source Gaussian-LIC2 codebase, with an empirical evaluation campaign
on physical multi-rig hardware in progress.

We believe the protocol design of Co-LIC2 opens a new axis for mobile
computing research: the joint optimization of perceptual data compression
and wireless transport for emerging 3D scene representations, extending
beyond SLAM to broader edge-assisted perception and digital-twin
applications.
